\begin{theorem}[径向二次观测的单样本离散可识别性：整数参数的支持严格分离]\label{thm:group-jg-radial-quadratic-single-sample-identifiability}
设 \(r>0\)。定义 Joukowsky--Gödel 椭圆映射
\[
J_r:\ \CC^\times\to\CC,\qquad J_r(z)=rz+r^{-1}z^{-1}.
\]
令 \(\Theta\sim \mathrm{Unif}[0,2\pi)\)，并置
\[
W_r:=J_r(e^{\mathrm{i}\Theta}),\qquad U_r:=|W_r|^2.
\]
则有恒等式
\[
U_r=r^2+r^{-2}+2\cos(2\Theta),
\qquad
\mathrm{supp}(U_r)=I_r:=\bigl[r^2+r^{-2}-2,\ \ r^2+r^{-2}+2\bigr].
\]
若进一步限制 \(r\in\ZZ_{\ge 2}\)，则族 \(\{I_r\}_{r\ge 2}\) 两两不交。更精确地，对每个整数 \(r\ge 2\) 令
\[
g_r:=\inf I_{r+1}-\sup I_r,
\]
则 \(g_r>0\) 且满足统一下界
\[
g_r\ge g_2=\frac{31}{36}.
\]
因此当 \(r\in\ZZ_{\ge 2}\) 时，单次无噪声观测 \(U_r=u\) 唯一确定该整数 \(r\)。
\end{theorem}

\begin{proof}
由定义
\[
W_r=r e^{\mathrm{i}\Theta}+r^{-1}e^{-\mathrm{i}\Theta},
\]
故
\[
|W_r|^2
=(r e^{\mathrm{i}\Theta}+r^{-1}e^{-\mathrm{i}\Theta})(r e^{-\mathrm{i}\Theta}+r^{-1}e^{\mathrm{i}\Theta})
=r^2+r^{-2}+e^{2\mathrm{i}\Theta}+e^{-2\mathrm{i}\Theta}
=r^2+r^{-2}+2\cos(2\Theta),
\]
从而 \(\mathrm{supp}(U_r)=I_r\) 成立。
\par
对整数 \(r\ge 2\)，计算
\[
g_r=\bigl((r+1)^2+(r+1)^{-2}-2\bigr)-\bigl(r^2+r^{-2}+2\bigr)
=2r-3+\bigl((r+1)^{-2}-r^{-2}\bigr)
=2r-3-\frac{2r+1}{r^2(r+1)^2}.
\]
令
\[
f(x):=\frac{2x+1}{x^2(x+1)^2}\quad(x>0).
\]
则
\[
\frac{f'(x)}{f(x)}
=\frac{2}{2x+1}-\frac{2}{x}-\frac{2}{x+1}
=\frac{2}{2x+1}\Bigl(1-\frac{(2x+1)^2}{x(x+1)}\Bigr)<0,
\]
因为 \((2x+1)^2-x(x+1)=3x^2+3x+1>0\)。故 \(f\) 在 \((0,\infty)\) 上严格递减，特别地对 \(r\ge 2\) 有
\[
f(r)\le f(2)=\frac{5}{36}.
\]
于是
\[
g_r=2r-3-f(r)\ge 2r-3-\frac{5}{36}
=\frac{72r-113}{36}\ge \frac{31}{36}.
\]
从而 \(g_r>0\)，即 \(I_r\cap I_{r+1}=\varnothing\)。又由于 \(r\mapsto r^2+r^{-2}\) 在 \([1,\infty)\) 上严格递增，区间端点随 \(r\) 单调右移，故全族两两不交。证毕。
\end{proof}

\begin{corollary}[有界噪声下的单样本恢复阈值]\label{cor:group-jg-radial-quadratic-bounded-noise-threshold}
沿用定理 \ref{thm:group-jg-radial-quadratic-single-sample-identifiability} 的记号。
设观测量带有加性噪声
\[
\widetilde U=U_r+\xi,\qquad |\xi|\le \sigma\ \ \text{几乎处处}.
\]
若 \(r\in\ZZ_{\ge 2}\) 且
\[
\sigma<\frac{1}{2}\inf_{r\ge 2}g_r=\frac{31}{72},
\]
则存在确定性解码规则使得单次观测 \(\widetilde U\) 即可无差错恢复 \(r\)。
\end{corollary}

\begin{proof}
由定理 \ref{thm:group-jg-radial-quadratic-single-sample-identifiability}，\(\inf_{r\ge 2}g_r=31/36\)。
若 \(\sigma<31/72\)，则相邻区间的 \(\sigma\)-膨胀
\[
I_r^{(\sigma)}:=[\inf I_r-\sigma,\ \sup I_r+\sigma]
\]
仍不相交：
\(
\sup I_r+\sigma<\inf I_{r+1}-\sigma
\)。
因此 \(\widetilde U\) 只能落入唯一一个膨胀区间中，解码唯一且必然正确。证毕。
\end{proof}

\begin{proposition}[离散寄存器的径向二次单样本全恢复]\label{prop:group-jg-radial-quadratic-prime-register-recovery}
设一类可交换投影算法的 Gödel 外置素数寄存器由某个正整数 \(N\ge 2\) 完全编码（即算法的全部离散信息可由 \(N\) 的素因子指数向量唯一恢复）。
将该整数作为椭圆参数 \(r=N\)，并按定理 \ref{thm:group-jg-radial-quadratic-single-sample-identifiability} 构造 \(U_N=|J_N(e^{\mathrm{i}\Theta})|^2\)。
则在无噪声情形下，单次观测 \(U_N=u\) 唯一确定 \(N\)，从而由唯一分解定理恢复寄存器指数向量并恢复算法本身。
若存在有界噪声且满足推论 \ref{cor:group-jg-radial-quadratic-bounded-noise-threshold} 的阈值，则上述结论仍成立。
\end{proposition}

\begin{proof}
由定理 \ref{thm:group-jg-radial-quadratic-single-sample-identifiability}，整数 \(N\ge 2\) 的支持区间 \(I_N\) 彼此不交，故单次观测 \(u\in I_N\) 唯一确定 \(N\)。
由整数唯一素因子分解，\(N\) 决定其素数寄存器指数向量；在假设下该指数向量决定算法。
噪声情形由推论 \ref{cor:group-jg-radial-quadratic-bounded-noise-threshold} 直接推出。证毕。
\end{proof}

\begin{theorem}[边界驱动的超二次收敛：连续参数的极值估计]\label{thm:group-jg-radial-quadratic-extreme-concentration}
仍在定理 \ref{thm:group-jg-radial-quadratic-single-sample-identifiability} 的设定下令
\[
c(r):=r^2+r^{-2},
\qquad U_r=c(r)+V,\qquad V:=2\cos(2\Theta).
\]
取 \(n\) 个独立同分布样本 \(U^{(1)},\dots,U^{(n)}\)，并定义
\[
\underline c_n:=\max_{1\le j\le n}U^{(j)}-2,\qquad
\overline c_n:=\min_{1\le j\le n}U^{(j)}+2.
\]
则必有 \(\underline c_n\le c(r)\le \overline c_n\)。
并且对任意 \(\varepsilon\in(0,2]\) 成立显式尾界
\[
\PP\bigl(c(r)-\underline c_n>\varepsilon\bigr)
=\Bigl(1-\frac{1}{\pi}\arccos\!\bigl(1-\varepsilon/2\bigr)\Bigr)^{n}
\le \exp\!\Bigl(-\frac{n}{\pi}\arccos\!\bigl(1-\varepsilon/2\bigr)\Bigr)
\le \exp\!\Bigl(-\frac{n}{\pi}\sqrt{\varepsilon}\Bigr),
\]
以及对称地
\[
\PP\bigl(\overline c_n-c(r)>\varepsilon\bigr)\le \exp\!\Bigl(-\frac{n}{\pi}\sqrt{\varepsilon}\Bigr).
\]
因此对任意 \(\eta\in(0,1)\)，取
\(
\varepsilon_\eta:=(\pi \log(2/\eta)/n)^2
\)
则
\[
\PP\bigl(|\overline c_n-\underline c_n|>2\varepsilon_\eta\bigr)\le \eta,
\]
从而在固定置信水平下 \(|\overline c_n-\underline c_n|\) 的典型尺度为 \(n^{-2}\)。
\end{theorem}

\begin{proof}
\(\underline c_n\le c(r)\le \overline c_n\) 由 \(V\in[-2,2]\) 立得。
设 \(M_n:=\max_j V^{(j)}\)，则
\[
c(r)-\underline c_n=c(r)-(\max_j(c(r)+V^{(j)})-2)=2-M_n.
\]
令 \(\Phi\sim\mathrm{Unif}[0,2\pi)\)。由于 \(2\Theta\bmod 2\pi\) 仍为均匀分布，有 \(V=2\cos\Phi\)。
对 \(\varepsilon\in(0,2]\)，
\[
\PP(V>2-\varepsilon)=\PP(\cos\Phi>1-\varepsilon/2)=\frac{1}{\pi}\arccos(1-\varepsilon/2),
\]
从而
\[
\PP(2-M_n>\varepsilon)=\PP(M_n\le 2-\varepsilon)
=\Bigl(1-\frac{1}{\pi}\arccos(1-\varepsilon/2)\Bigr)^{n}.
\]
不等式 \((1-x)^n\le e^{-nx}\) 给出第一条指数上界。
此外，由 \(\cos t\ge 1-t^2/2\)（\(t\in[0,\pi]\)）可知当 \(\cos t=1-\varepsilon/2\) 时 \(t\ge \sqrt{\varepsilon}\)，即
\(
\arccos(1-\varepsilon/2)\ge \sqrt{\varepsilon}
\)，从而得到最后一条上界。
对 \(\overline c_n\) 同理。
最后，事件
\(
\{c(r)-\underline c_n\le \varepsilon_\eta\}\cap\{\overline c_n-c(r)\le \varepsilon_\eta\}
\)
蕴含 \(|\overline c_n-\underline c_n|\le 2\varepsilon_\eta\)，并由并集界得到概率陈述。证毕。
\end{proof}

\begin{proposition}[分布支持分离导致的 \texorpdfstring{$\mathrm{KL}$}{KL} 发散]\label{prop:group-jg-radial-quadratic-kl-infinite}
记 \(\mathsf P_r\) 为随机变量 \(U_r\) 的分布。则对任意 \(r\neq s\) 有
\[
D_{\mathrm{KL}}(\mathsf P_r\Vert \mathsf P_s)=+\infty,
\qquad
D_{\mathrm{KL}}(\mathsf P_s\Vert \mathsf P_r)=+\infty.
\]
特别地，只要 \(c(r)\neq c(s)\)，则 \(\mathsf P_r\) 与 \(\mathsf P_s\) 的支持区间不同，故两分布在 Lebesgue 分解意义下互不绝对连续。
\end{proposition}

\begin{proof}
由定理 \ref{thm:group-jg-radial-quadratic-single-sample-identifiability}，\(\mathrm{supp}(U_r)=I_r=[c(r)-2,c(r)+2]\)。
若 \(r\neq s\)，则 \(c(r)\neq c(s)\)（因 \(x\mapsto x^2+x^{-2}\) 在 \([1,\infty)\) 上严格递增），从而 \(I_r\neq I_s\)。
于是存在可测集 \(A\subset I_r\setminus I_s\) 使 \(\mathsf P_r(A)>0\) 而 \(\mathsf P_s(A)=0\)。
此时 \(\mathsf P_r\not\ll \mathsf P_s\)，故 \(D_{\mathrm{KL}}(\mathsf P_r\Vert \mathsf P_s)=+\infty\)；交换 \(r,s\) 同理。证毕。
\end{proof}

\begin{conjecture}[非交换投影的低阶可观测解码泛化]\label{conj:group-jg-noncommutative-low-order-observable-decoding}
设 \(\mathcal A\) 为由素数寄存器 Gödel 编码生成的非交换投影算法类，其对应的单位圆信息经一族可计算的代数映射 \(F_{\alpha}\)（参数 \(\alpha\) 取 Gödel 整数或其有限素因子指数向量）投影为代数曲线族 \(\Gamma_\alpha\subset\CC\)。
猜想存在一个与 \(\alpha\) 无关的固定次数多项式可观测量 \(\mathcal O(w,\bar w)\in\RR[w,\bar w]\)（例如次数 \(\le 4\)），使得在离散参数子集上，\(\mathcal O\) 的像分布族呈现支持严格分离，从而产生常数样本量的算法全恢复机制，与命题 \ref{prop:group-jg-radial-quadratic-prime-register-recovery} 的可交换情形一致。
\end{conjecture}

\begin{remark}[可复现核验脚本]\label{rem:group-jg-radial-quadratic-identifiability-audit-script}
定理 \ref{thm:group-jg-radial-quadratic-single-sample-identifiability}--命题 \ref{prop:group-jg-radial-quadratic-kl-infinite}
中的区间分离常数、噪声阈值、极值尾界与支持分离的有限窗核验由脚本
\texttt{scripts/exp\_group\_jg\_radial\_quadratic\_identifiability\_audit.py}
给出，对应证书文件为
\texttt{artifacts/export/group\_jg\_radial\_quadratic\_identifiability\_audit.json}。
\end{remark}

